% Source - https://tex.stackexchange.com/a/457624
% Posted by user121799
% Retrieved 2026-03-17, License - CC BY-SA 4.0

\documentclass[tikz,border=8mm]{standalone}
\usetikzlibrary{shapes.geometric,shapes.symbols,fit,positioning,shadows}
% https://tex.stackexchange.com/a/12039/121799
\makeatletter
% https://tex.stackexchange.com/a/103691/121799
\pgfdeclareshape{document}{
\inheritsavedanchors[from=rectangle] % this is nearly a rectangle
\inheritanchorborder[from=rectangle]
\inheritanchor[from=rectangle]{center}
\inheritanchor[from=rectangle]{north}
\inheritanchor[from=rectangle]{north east}
\inheritanchor[from=rectangle]{north west}
\inheritanchor[from=rectangle]{south}
\inheritanchor[from=rectangle]{south east}
\inheritanchor[from=rectangle]{south west}
\inheritanchor[from=rectangle]{west}
\inheritanchor[from=rectangle]{east}
\backgroundpath{%
\southwest \pgf@xa=\pgf@x \pgf@ya=\pgf@y
\northeast \pgf@xb=\pgf@x \pgf@yb=\pgf@y
\pgf@xc=\pgf@xb \advance\pgf@xc by-5pt % this should be a parameter
\pgf@yc=\pgf@ya \advance\pgf@yc by5pt
\pgfpathmoveto{\pgfpoint{\pgf@xa}{\pgf@ya}}
\pgfpathlineto{\pgfpoint{\pgf@xa}{\pgf@yb}}
\pgfpathlineto{\pgfpoint{\pgf@xb}{\pgf@yb}}
\pgfpathlineto{\pgfpoint{\pgf@xb}{\pgf@yc}}
\pgfpathlineto{\pgfpoint{\pgf@xc}{\pgf@ya}}
\pgfpathclose
% add little corner
\pgfpathmoveto{\pgfpoint{\pgf@xc}{\pgf@ya}}
\pgfpathlineto{\pgfpoint{\pgf@xc}{\pgf@yc}}
\pgfpathlineto{\pgfpoint{\pgf@xb}{\pgf@yc}}
\pgfpathclose
}
}
\makeatother

\begin{document}

\tikzset{doc/.style={document,fill=blue!10,draw,thin,minimum
height=1.2cm,align=center},
pics/.cd,
pack/.style={code={%
\draw[fill=blue!50,opacity=0.2] (0,0) -- (0.5,-0.25) -- (0.5,0.25) -- (0,0.5) -- cycle;
\draw[fill=blue!50,opacity=0.2] (0,0) -- (-0.5,-0.25) -- (-0.5,0.25) -- (0,0.5) -- cycle;
\draw[fill=blue!60,opacity=0.2] (0,0) -- (-0.5,-0.25) -- (0,-0.5) -- (0.5,-0.25) -- cycle;
\draw[fill=blue!60] (0,0) -- (0.25,0.125) -- (0,0.25) -- (-0.25,0.125) -- cycle;
\draw[fill=blue!50] (0,0) -- (0.25,0.125) -- (0.25,-0.125) -- (0,-0.25) -- cycle;
\draw[fill=blue!50] (0,0) -- (-0.25,0.125) -- (-0.25,-0.125) -- (0,-0.25) -- cycle;
\draw[fill=blue!50,opacity=0.2] (0,-0.5) -- (0.5,-0.25) -- (0.5,0.25) -- (0,0) -- cycle;
 \draw[fill=blue!50,opacity=0.2] (0,-0.5) -- (-0.5,-0.25) -- (-0.5,0.25) -- (0,0) -- cycle;
\draw[fill=blue!60,opacity=0.2] (0,0.5) -- (-0.5,0.25) -- (0,0) -- (0.5,0.25) -- cycle;
}}}

\begin{tikzpicture}[
    font=\sffamily,
    every label/.append style={font=\small\sffamily,align=center}]

% I/O layer nodes
\node[doc,fill=red!10]
    (GPIO)
    {GPIO\\Interface\\(Mutex)};
\node[right=2cm of GPIO,doc,fill=blue!10]
    (PWM)
    {PWM\\Interface};
\node[right=2cm of PWM,doc,fill=blue!10]
    (Console)
    {Console};

% I/O thread layer nodes
\node[below left=3cm and 2cm of GPIO,doc,fill=blue!10]
    (BeamThread)
    {Beam\\Thread};
\node[right=2cm of BeamThread,doc,fill=blue!10]
    (DistanceThread)
    {Distance\\Thread};
\node[right=2cm of DistanceThread,doc,fill=blue!10]
    (ConsoleThread)
    {Console\\Thread};
\node[right=2cm of ConsoleThread,doc,fill=blue!10]
    (DisplayThread)
    {Display\\Thread};
\node[right=2cm of DisplayThread,doc,fill=blue!10]
    (LoggingThread)
    {Logging\\Thread};

% State layer nodes
\node[below right=3cm and 2cm of BeamThread,doc,fill=red!10]
    (DistanceState)
    {Distance\\State\\(Mutex)};
\node[right=2cm of DistanceState,doc,fill=red!10]
    (LoggingState)
    {Logging\\State\\(Mutex)};
\node[right=2cm of LoggingState,doc,fill=red!10]
    (SafeState)
    {Safe\\State\\(Mutex)};


% \draw[-latex] (ClientBundle) -- (HTML) node[midway,below,font=\small\sffamily]{Hydrate};
% \draw (SensorThread) -- (ConsoleThread);

% links between layers 1 and 2
\draw[latex-] (ConsoleThread) -- (Console);
\draw[latex-] (BeamThread) -- (GPIO);
\draw[-latex] (BeamThread) -- (PWM);
\draw[latex-latex] (DistanceThread) -- (GPIO);
\draw[-latex] (DisplayThread) -- (GPIO);
\draw[-latex] (LoggingThread) -- (Console);

% links between layers 2 and 3
\draw[latex-] (BeamThread) -- (DistanceState);
\draw[-latex] (BeamThread) -- (LoggingState);
\draw[latex-latex] (BeamThread) -- (SafeState);

\draw[-latex] (DistanceThread) -- (DistanceState);
\draw[-latex] (DistanceThread) -- (LoggingState);
\draw[-latex] (DistanceThread) -- (SafeState);

\draw[-latex] (ConsoleThread) -- (LoggingState);
\draw[-latex] (ConsoleThread) -- (SafeState);

\draw[latex-] (DisplayThread) -- (DistanceState);
\draw[-latex] (DisplayThread) -- (LoggingState);
\draw[latex-] (DisplayThread) -- (SafeState);

\draw[latex-] (LoggingThread) -- (LoggingState);

\node[draw,dashed,rounded corners,
    fit=(GPIO)(PWM)(Console),
    inner sep=10pt,label={left:{I/O\\Interfaces}}]
    (fit1)
    {};
\node[draw,dashed,rounded corners,
    fit=(BeamThread)(DistanceThread)(ConsoleThread),
    inner sep=10pt,label={left:{Input\\Threads}}]
    (fit2)
    {};
\node[draw,dashed,rounded corners,
    fit=(DisplayThread)(LoggingThread),
    inner sep=10pt,label={left:{Output\\Threads}}]
    (fit3)
    {};
\node[draw,dashed,rounded corners,
    fit=(DistanceState)(LoggingState)(SafeState),
    inner sep=10pt,label={left:{Shared\\State}}]
    (fit4)
    {};

\end{tikzpicture}
\end{document}
